\documentclass[11pt,a4paper]{article}

% ============================================
% PACCHETTI
% ============================================
\usepackage[utf8]{inputenc}
\usepackage[T1]{fontenc}
\usepackage[italian]{babel}
\usepackage{geometry}
\usepackage{hyperref}
\usepackage{xcolor}
\usepackage{listings}
\usepackage{tikz}
\usepackage{booktabs}
\usepackage{longtable}
\usepackage{array}
\usepackage{fancyhdr}
\usepackage{titlesec}
\usepackage{amsmath}
\usepackage{graphicx}
\usepackage{float}
\usepackage{enumitem}

% ============================================
% CONFIGURAZIONE PAGINA
% ============================================
\geometry{margin=2.5cm}
\hypersetup{
    colorlinks=true,
    linkcolor=blue,
    filecolor=magenta,      
    urlcolor=cyan,
    pdftitle={Deck/Flashcard Component - Technical Documentation},
    pdfauthor={Silvi Development Team}
}

% ============================================
% COLORI
% ============================================
\definecolor{codegreen}{rgb}{0,0.6,0}
\definecolor{codegray}{rgb}{0.5,0.5,0.5}
\definecolor{codepurple}{rgb}{0.58,0,0.82}
\definecolor{backcolour}{rgb}{0.95,0.95,0.92}
\definecolor{darkblue}{rgb}{0.0,0.0,0.6}
\definecolor{deepred}{rgb}{0.6,0,0}
\definecolor{deepgreen}{rgb}{0,0.5,0}

% ============================================
% STILE LISTATI CODICE
% ============================================
\lstdefinestyle{mystyle}{
    backgroundcolor=\color{backcolour},   
    commentstyle=\color{codegreen},
    keywordstyle=\color{magenta},
    numberstyle=\tiny\color{codegray},
    stringstyle=\color{codepurple},
    basicstyle=\ttfamily\footnotesize,
    breakatwhitespace=false,         
    breaklines=true,                 
    captionpos=b,                    
    keepspaces=true,                 
    numbers=left,                    
    numbersep=5pt,                  
    showspaces=false,                
    showstringspaces=false,
    showtabs=false,                  
    tabsize=2
}
\lstset{style=mystyle}

% ============================================
% TIKZ LIBRERIE
% ============================================
\usetikzlibrary{shapes.geometric, arrows, positioning, fit, backgrounds, calc}

% ============================================
% STILE TITOLI
% ============================================
\titleformat{\section}{\Large\bfseries\color{darkblue}}{\thesection}{1em}{}
\titleformat{\subsection}{\large\bfseries\color{darkblue}}{\thesubsection}{1em}{}
\titleformat{\subsubsection}{\normalsize\bfseries\color{darkblue}}{\thesubsubsection}{1em}{}

% ============================================
% HEADER E FOOTER
% ============================================
\pagestyle{fancy}
\fancyhf{}
\fancyhead[L]{\leftmark}
\fancyhead[R]{Deck/Flashcard Component}
\fancyfoot[C]{\thepage}
\renewcommand{\headrulewidth}{0.4pt}
\renewcommand{\footrulewidth}{0.4pt}

% ============================================
% TITOLO
% ============================================
\title{
    \Huge\textbf{Deck/Flashcard Component}\\[0.5em]
    \Large\textit{Documentazione Tecnica Completa}\\[1em]
    \large Silvi - Spaced Repetition System
}
\author{
    \textbf{Extra Tracker Development Team}\\
    \texttt{Versione 1.0 - Febbraio 2026}
}
\date{\today}

% ============================================
% INIZIO DOCUMENTO
% ============================================
\begin{document}

\maketitle
\thispagestyle{empty}

\vfill

\begin{abstract}
\noindent
Questo documento fornisce una documentazione tecnica esaustiva del componente \textbf{Deck/Flashcard} del sistema Silvi. Il componente implementa un sistema completo di \textit{spaced repetition} con supporto per algoritmi SM-2, FSRS e Leitner, generazione AI di flashcard da PDF, e multi-modalita' di studio (flashcard, quiz, typing, mix). Il documento copre l'architettura completa sia backend (Node.js/Express/MongoDB) che frontend (React/TypeScript), i flussi di comunicazione, gli algoritmi core, e le integrazioni con servizi AI.
\end{abstract}

\vfill
\newpage

% ============================================
% INDICE
% ============================================
\tableofcontents
\newpage

% ============================================
% SEZIONE 1: INTRODUZIONE
% ============================================
\section{Introduzione}\label{sec:intro}

\subsection{Panoramica del Componente}

Il componente \textbf{Deck/Flashcard} e' il nucleo del sistema di apprendimento di Silvi. Implementa un sistema completo di gestione mazzi di flashcard con le seguenti caratteristiche principali:

\begin{itemize}[leftmargin=*]
    \item \textbf{Spaced Repetition}: Implementazione degli algoritmi SM-2, FSRS e Leitner
    \item \textbf{Multi-Modalita'}: Flashcard, Quiz (multiple choice), Typing, Mix, Sprint, Focus, Exam
    \item \textbf{AI Integration}: Generazione automatica di flashcard da PDF tramite OpenAI
    \item \textbf{Source Grounding}: Tracciamento della fonte originale nel PDF per ogni flashcard
    \item \textbf{Exam Integration}: Collegamento diretto con il sistema di gestione esami
    \item \textbf{Multi-Tenancy}: Isolamento completo dei dati per utente
\end{itemize}

\subsection{Architettura Generale}

Il sistema segue un'architettura \textbf{Client-Server} a tre livelli:

\begin{center}
\begin{tikzpicture}[
    box/.style={rectangle, rounded corners, draw=darkblue, fill=blue!5, minimum width=3cm, minimum height=1cm, align=center},
    arrow/.style={->, >=stealth, thick}
]
    \node[box] (client) at (0,0) {\textbf{Client}\\React + TypeScript};
    \node[box] (api) at (5,0) {\textbf{API Layer}\\Express + Middleware};
    \node[box] (service) at (10,0) {\textbf{Service Layer}\\Business Logic};
    \node[box] (db) at (10,-2.5) {\textbf{Database}\\MongoDB + Mongoose};
    \node[box] (ai) at (5,-2.5) {\textbf{AI Services}\\OpenAI};
    
    \draw[arrow] (client) -- (api) node[midway,above] {\small HTTPS/JSON};
    \draw[arrow] (api) -- (service) node[midway,above] {\small tenantScope};
    \draw[arrow] (service) -- (db) node[midway,right] {\small Queries};
    \draw[arrow, dashed] (service) -- (ai) node[midway,above] {\small PDF/Gen};
\end{tikzpicture}
\end{center}

\subsection{Struttura dei File}

\begin{lstlisting}[language=bash, caption={Struttura file componente Deck/Flashcard}]
# BACKEND
server/
|-- models/
|   |-- Deck.js              # Schema MongoDB (cards embedded)
|   |-- Exam.js              # Schema esami (relazione opzionale)
|-- controllers/
|   |-- studyController.js   # HTTP request handlers
|-- services/
|   |-- studyService.js      # Business logic + AI generation
|-- routes/
|   |-- study.js             # Route definitions (/api/study)
|-- plugins/
|   |-- multiTenancy.js      # Tenant isolation plugin

# FRONTEND
src/features/study/
|-- services/
|   |-- studyService.ts      # API client + normalization
|-- components/
|   |-- DeckCard/            # Visualizzazione deck
|   |-- Flashcard/           # Studio flashcard
|   |-- DeckSections/        # Organizzazione folder
|-- pages/
|   |-- DeckDetailPage.tsx   # Pagina dettaglio deck
|   |-- StudySessionPage.tsx # Sessione di studio
\end{lstlisting}

% ============================================
% SEZIONE 2: BACKEND
% ============================================
\section{Backend Architecture}\label{sec:backend}

\subsection{Data Model: Deck Schema}\label{subsec:deck-model}

Il modello Deck utilizza uno schema MongoDB con \textbf{embedded documents} per le cards. Questa scelta progettuale e' ottimale perche':
\begin{itemize}
    \item Le cards sono sempre accedute nel contesto del loro deck
    \item Non ci sono relazioni many-to-many complesse
    \item Le operazioni di lettura sono atomiche (single document)
\end{itemize}

\begin{lstlisting}[language=JavaScript, caption={Schema Deck - Card Embedded}]
const cardSchema = new mongoose.Schema({
    // Contenuto
    front: { type: String, required: true, maxlength: 2000 },
    back:  { type: String, required: true, maxlength: 4000 },
    
    // Parametri SM-2
    easinessFactor: { type: Number, default: 2.5, min: 1.3 },
    interval:       { type: Number, default: 0 },
    repetitions:    { type: Number, default: 0 },
    nextReviewDate: { type: Date, default: Date.now },
    status: { type: String, enum: ['new','learning','review','mastered'] },
    
    // Parametri FSRS
    stability: { type: Number, default: 0.4 },
    
    // Parametri Leitner
    box: { type: Number, default: 1, min: 1 },
    
    // History tracking
    lastReviewed: Date,
    reviewHistory: [{
        date: Date,
        rating: { type: Number, min: 1, max: 5 },
        interval: Number,
        easinessFactor: Number,
        repetitions: Number,
        algorithm: String
    }],
    
    // Source Grounding (tracciamento fonte PDF)
    sourceMetadata: {
        pageNumber: { type: Number, required: true, min: 1 },
        originalText: { type: String, required: true, minlength: 150 }
    }
}, { _id: true });
\end{lstlisting}

\textbf{Deck Schema (parent)}:

\begin{lstlisting}[language=JavaScript]
const deckSchema = new mongoose.Schema({
    examId:     { type: ObjectId, ref: 'Exam', default: null, index: true },
    folderId:   { type: ObjectId, ref: 'Folder', default: null, index: true },
    title:      { type: String, required: true, maxlength: 120 },
    description:{ type: String, maxlength: 500 },
    pdfUrl:     { type: String, default: '' },
    extractedText: { type: String, default: '', select: false },
    tags:       { type: [String], default: [], validate: v => v.length <= 5 },
    cards:      { type: [cardSchema], default: [] },
    
    // Algoritmo e AI Settings
    algorithm:  { type: String, enum: ['sm2','fsrs','leitner','anki'], 
                  default: 'sm2' },
    aiSettings: {
        style: { type: String, enum: ['comprehensive','conceptual','factual','application'] },
        difficulty: { type: String, enum: ['easy','medium','hard','mixed'] },
        questionTypes: { type: [String], default: ['definition','concept','relationship'] }
    }
}, { timestamps: true });
\end{lstlisting}

\subsection{Multi-Tenancy Implementation}\label{subsec:multitenancy}

Il sistema implementa il multi-tenancy tramite \textbf{Row-Level Security}: ogni query include automaticamente il filtro user.

\begin{lstlisting}[language=JavaScript, caption={Multi-Tenancy Plugin}]
function multiTenancyPlugin(schema, options = {}) {
    const tenantField = options.tenantField || 'user';
    
    // Aggiungi campo user a tutti i documenti
    schema.add({
        [tenantField]: {
            type: mongoose.Schema.Types.ObjectId,
            ref: 'User',
            required: options.required,
            index: true
        }
    });
    
    // Pre-hook per auto-filtrare per tenant
    schema.pre(['find', 'findOne', 'findOneAndUpdate', 'countDocuments'], 
        function() {
            const tenantScope = this.options?.tenantScope;
            if (tenantScope?.userId) {
                this.where({ [tenantField]: tenantScope.userId });
            }
        }
    );
}
\end{lstlisting}

\textbf{Tenant Context Middleware}:
\begin{lstlisting}[language=JavaScript]
function tenantContext(options = {}) {
    return (req, res, next) => {
        // Estrai userId dal JWT verificato da requireAuth
        const userId = req.user?.id || req.user?._id;
        
        if (options.required && !userId) {
            return res.status(401).json({ 
                success: false, 
                error: { message: 'Authentication required' } 
            });
        }
        
        // Espandi req con tenantScope
        req.tenantScope = { userId };
        next();
    };
}
\end{lstlisting}

\textbf{Utilizzo nelle routes}:
\begin{lstlisting}[language=JavaScript]
const { requireAuth } = require('../middleware/auth');
const { tenantContext } = require('../middleware/tenantContext');

router.use(requireAuth);
router.use(tenantContext({ required: true }));

router.get('/dashboard', studyController.getDashboard);
// Ogni query nel controller avra' automaticamente filter: { user: userId }
\end{lstlisting}

\subsection{Service Layer: StudyService}\label{subsec:study-service}

Il StudyService estende BaseService e implementa la business logic completa.

\begin{lstlisting}[language=JavaScript, caption={StudyService - Architettura}]
class StudyService extends BaseService {
    constructor() {
        super(Deck, {
            searchFields: ['title', 'description', 'tags'],
            defaultSort: { createdAt: -1 },
            entityName: 'Mazzo',
        });
    }
    
    // ================================
    // CRUD Deck
    // ================================
    async createDeck(tenantScope, { examId, title, description, tags }) {
        // Validazione ownership esame se specificato
        if (examId) {
            await this._validateExamOwnership(tenantScope, examId);
        }
        
        return this.create(tenantScope, {
            examId: examId || null,
            title,
            description,
            tags,
        });
    }
    
    async updateDeck(tenantScope, deckId, updates) {
        const userId = this._getUserId(tenantScope);
        const deck = await Deck.findOne({ _id: deckId, user: userId });
        
        if (!deck) throw AppError.notFound('Mazzo');
        
        // Gestione aggiornamento folderId
        if (updates.folderId !== undefined) {
            if (updates.folderId !== null && updates.folderId !== '') {
                const Folder = require('../models/Folder');
                const folder = await Folder.findOne({ 
                    _id: updates.folderId, 
                    user: userId 
                });
                if (!folder) throw AppError.notFound('Cartella');
                deck.folderId = updates.folderId;
            } else {
                deck.folderId = null;
            }
        }
        
        // Gestione aggiornamento examId
        if (updates.examId !== undefined) {
            if (updates.examId !== null && updates.examId !== '') {
                const exam = await Exam.findOne({ 
                    _id: updates.examId, 
                    user: userId 
                });
                if (!exam) throw AppError.notFound('Esame');
                deck.examId = updates.examId;
            } else {
                deck.examId = null;
            }
        }
        
        await deck.save();
        return this._serializeDeck(deck);
    }
}
\end{lstlisting}

\subsection{Spaced Repetition Algorithms}\label{subsec:algorithms}

\subsubsection{Algoritmo SM-2 (Default)}

L'algoritmo SM-2 e' implementato nella AlgorithmFactory:

\begin{lstlisting}[language=JavaScript, caption={SM-2 Algorithm Implementation}]
class SM2Algorithm {
    processReview(card, quality) {
        const MIN_EF = 1.3;
        let { easinessFactor, interval, repetitions } = card;
        
        // Aggiorna easiness factor secondo formula SM-2
        easinessFactor = Math.max(
            MIN_EF,
            easinessFactor + (0.1 - (5 - quality) * (0.08 + (5 - quality) * 0.02))
        );
        
        // Calcola nuovo intervallo
        if (quality < 3) {
            // Risposta difficile: resetta
            repetitions = 0;
            interval = 1;
        } else {
            // Risposta corretta: incrementa intervallo
            repetitions += 1;
            if (repetitions === 1) interval = 1;
            else if (repetitions === 2) interval = 6;
            else interval = Math.round(interval * easinessFactor);
        }
        
        // Calcola prossima data revisione
        const nextReviewDate = new Date();
        nextReviewDate.setDate(nextReviewDate.getDate() + interval);
        
        return {
            easinessFactor,
            interval,
            repetitions,
            nextReviewDate,
            status: this._resolveStatus(quality, repetitions, interval)
        };
    }
}
\end{lstlisting}

\textbf{Formula SM-2}:
\begin{align}
EF' &= EF + (0.1 - (5 - q) \times (0.08 + (5 - q) \times 0.02)) \\
EF' &\geq 1.3 \\
I_n &= \begin{cases} 
1 & \text{if } n = 1 \\
6 & \text{if } n = 2 \\
I_{n-1} \times EF & \text{if } n > 2 \text{ and } q \geq 3 \\
1 & \text{if } q < 3
\end{cases}
\end{align}

Dove $q$ e' la qualita' della risposta (1-5) e $n$ e' il numero di ripetizioni.

\subsubsection{Selezione Cards per Sessione}

\begin{lstlisting}[language=JavaScript, caption={Session Card Selection}]
_selectSessionCards({ cards, dueCards, mode, focus, limit, now }) {
    // Cards effettivamente in scadenza
    const due = cards.filter(c => new Date(c.nextReviewDate) <= now);
    
    switch (focus) {
        case 'due':
            // Solo cards in scadenza
            return due.slice(0, limit);
            
        case 'weak':
            // Cards con easiness factor basso (difficili)
            return cards
                .filter(c => c.easinessFactor < 2.0)
                .sort((a, b) => a.easinessFactor - b.easinessFactor)
                .slice(0, limit);
                
        case 'smart':
        default:
            // Priorita: due > new > weak
            const newCards = cards.filter(c => c.status === 'new');
            const weakCards = cards.filter(c => 
                c.easinessFactor < 2.0 && !due.includes(c)
            );
            return [...due, ...newCards, ...weakCards].slice(0, limit);
    }
}
\end{lstlisting}

\subsection{AI Integration: Magic Generate}\label{subsec:magic-generate}

Il sistema genera flashcard da PDF usando OpenAI con una pipeline multi-stage:

\begin{lstlisting}[language=JavaScript, caption={Magic Generate V3 - Pipeline Completa}]
async generateCardsFromPDF(tenantScope, deckId, pdfFilePath) {
    const userId = this._getUserId(tenantScope);
    
    // 1. Verifica ownership
    const deck = await Deck.findOne({ _id: deckId, user: userId });
    if (!deck) throw AppError.notFound('Mazzo');
    
    // 2. Estrai testo dal PDF
    const pdfBuffer = await fs.readFile(pdfFilePath);
    const pdfData = await pdfCacheService.parsePDF(pdfFilePath, pdfBuffer);
    const pdfText = this._formatPdfTextWithPages(pdfData);
    
    if (!pdfText || pdfText.trim().length < 100) {
        throw AppError.validation('PDF non contiene abbastanza testo');
    }
    
    // 3. Analisi strutturale (Blueprint)
    const blueprint = await this._analyzeDocumentStructure(pdfText);
    
    // 4. Semantic Chunking (12k caratteri per chunk)
    const semanticChunks = this._createSemanticChunks(pdfText);
    
    // 5. Estrazione concetti locale (senza AI)
    const globalConcepts = this._extractConceptsLocally(pdfText);
    
    // 6. Batch Generation (2 chunk per chiamata API)
    const batches = this._createBatches(semanticChunks, BATCH_SIZE = 2);
    let allGeneratedCards = [];
    const usedConcepts = new Set(globalConcepts.slice(0, 10));
    
    for (const batch of batches) {
        const targetCards = this._calculateBatchTarget(batch);
        
        const batchCards = await this._generateCardsBatch(
            batch, blueprint, targetCards, usedConcepts
        );
        
        // Traccia concetti usati per evitare duplicati
        batchCards.forEach(card => {
            const conceptKey = this._extractConceptKey(card.front);
            if (conceptKey) usedConcepts.add(conceptKey);
        });
        
        allGeneratedCards.push(...batchCards);
    }
    
    // 7. Deduplica semantica (Jaccard similarity < 0.50)
    const uniqueCards = this._deduplicateCards(allGeneratedCards);
    
    // 8. Validazione qualita' e salvataggio
    const validCards = uniqueCards
        .filter(card => this._validateCardQuality(card))
        .slice(0, MAX_TOTAL_CARDS)
        .map(card => ({
            front: card.front.trim(),
            back: card.back.trim(),
            status: 'new',
            nextReviewDate: new Date(),
            easinessFactor: DEFAULT_EASINESS_FACTOR,
            interval: 0,
            repetitions: 0,
        }));
    
    deck.cards.push(...validCards);
    await deck.save();
    
    return { generatedCount: validCards.length, deck };
}
\end{lstlisting}

\textbf{Prompt Engineering}:
\begin{lstlisting}[language=bash]
SYSTEM PROMPT:
"Sei un esperto di didattica. Crea flashcard di alta qualita' dal testo 
fornito. Regole:
- Domande specifiche e non ambigue
- Risposte concise ma complete  
- Usa LaTeX per formule matematiche
- Una domanda per concetto"

USER PROMPT:
"Crea {targetCards} flashcard da questa sezione:
---
{chunk.text}
---

Concetti gia' usati: {Array.from(usedConcepts).join(', ')}
Tipo: {selectedQuestionType}

Rispondi in JSON: { "flashcards": [{"front":"...","back":"..."}] }"
\end{lstlisting}

\subsection{API Endpoints}\label{subsec:api-endpoints}

\begin{longtable}{|p{4cm}|p{1.5cm}|p{8cm}|}
\hline
\textbf{Endpoint} & \textbf{Method} & \textbf{Descrizione} \\
\hline
\endfirsthead
\hline
\textbf{Endpoint} & \textbf{Method} & \textbf{Descrizione} \\
\hline
\endhead
\texttt{/api/study/dashboard} & GET & Recupera tutti i deck con conteggio cards in scadenza \\
\hline
\texttt{/api/study/:id} & GET & Recupera singolo deck con tutte le cards \\
\hline
\texttt{/api/study} & POST & Crea nuovo deck \\
\hline
\texttt{/api/study/:id} & PATCH & Aggiorna deck (title, desc, tags, folderId, examId) \\
\hline
\texttt{/api/study/:id} & DELETE & Elimina deck \\
\hline
\texttt{/api/study/:id/session} & GET & Recupera sessione di studio \\
\hline
\texttt{/api/study/:id/cards} & POST & Aggiunge card al deck \\
\hline
\texttt{/api/study/:id/cards/:cardId} & PUT & Modifica card esistente \\
\hline
\texttt{/api/study/:id/cards/:cardId} & DELETE & Elimina card \\
\hline
\texttt{/api/study/:id/cards/reorder} & PUT & Riordina cards (array cardIds) \\
\hline
\texttt{/api/study/:id/cards/insert} & POST & Inserisce card in posizione specifica \\
\hline
\texttt{/api/study/:id/review} & POST & Processa review SM-2 (cardId, rating 1-5) \\
\hline
\texttt{/api/study/:id/verify-answer} & POST & Verifica risposta typing mode \\
\hline
\texttt{/api/study/:id/generate-pdf} & POST & Upload PDF + generazione AI \\
\hline
\texttt{/api/study/:id/chat} & POST & Chat con AI Tutor (RAG sul PDF) \\
\hline
\end{longtable}

% ============================================
% SEZIONE 3: FRONTEND
% ============================================
\section{Frontend Architecture}\label{sec:frontend}

\subsection{Service Layer: studyService.ts}\label{subsec:frontend-service}

Il frontend utilizza un Service Pattern con normalizzazione difensiva dei dati.

\begin{lstlisting}[language=TypeScript, caption={StudyService - Frontend}]
export interface Card {
    id: string;
    front: string;
    back: string;
    options?: string[];  // Per quiz mode
    easinessFactor: number;
    interval: number;
    repetitions: number;
    nextReviewDate: string;
    status: CardStatus;
    sourceMetadata?: {
        pageNumber: number;
        originalText: string;
    };
}

export interface Deck {
    id: string;
    examId?: string;
    folderId?: string | null;
    title: string;
    description?: string;
    pdfUrl?: string | null;
    tags: string[];
    cards: Card[];
    totalCards: number;
    dueCount: number;
    algorithm?: string;
}

class StudyService {
    private baseUrl = '/study';
    
    // Normalizzazione difensiva
    private normalizeCard(raw: any): Card {
        let sourceMetadata: Card['sourceMetadata'] = undefined;
        
        // Supporta sia camelCase che snake_case
        const sourceMeta = raw.sourceMetadata || raw.source_metadata;
        if (sourceMeta && typeof sourceMeta === 'object') {
            const pageNumber = Number.isFinite(Number(
                sourceMeta.pageNumber ?? sourceMeta.page_number
            )) ? Number(sourceMeta.pageNumber ?? sourceMeta.page_number) : undefined;
            
            const originalText = typeof (
                sourceMeta.originalText ?? sourceMeta.original_text
            ) === 'string' ? (sourceMeta.originalText ?? sourceMeta.original_text).trim() : undefined;
            
            if (pageNumber && pageNumber > 0 && originalText && originalText.length >= 20) {
                sourceMetadata = { pageNumber, originalText };
            }
        }
        
        return {
            id: raw.id || raw._id,
            front: raw.front || '',
            back: raw.back || '',
            options: Array.isArray(raw.options) ? raw.options : undefined,
            easinessFactor: safeNumber(raw.easinessFactor, 2.5),
            interval: safeNumber(raw.interval, 0),
            repetitions: safeNumber(raw.repetitions, 0),
            nextReviewDate: raw.nextReviewDate || new Date().toISOString(),
            status: raw.status || 'new',
            sourceMetadata,
        };
    }
    
    async getSession(deckId: string, options: SessionOptions): Promise<StudySession> {
        const params = new URLSearchParams();
        if (options.mode) params.set('mode', options.mode);
        if (options.focus) params.set('focus', options.focus);
        if (options.limit) params.set('limit', String(options.limit));
        
        const response = await apiClient.get(
            `${this.baseUrl}/${deckId}/session?${params}`
        );
        return normalizeSession(unwrap(response, 'Errore recupero sessione'));
    }
    
    async submitReview(deckId: string, payload: ReviewPayload): Promise<ReviewResult> {
        const response = await apiClient.post(
            `${this.baseUrl}/${deckId}/review`, 
            payload
        );
        return unwrap(response, 'Errore salvataggio review');
    }
}
\end{lstlisting}

\subsection{API Client con Mutex Pattern}\label{subsec:api-client}

Il client API implementa il Mutex Pattern per gestire il refresh token.

\begin{lstlisting}[language=TypeScript, caption={Mutex Pattern per Refresh Token}]
// Stato del mutex
let isRefreshing = false;
let isSessionDead = false;
let failedQueue: Array<{ resolve: Function; reject: Function }> = [];

// Processa la coda delle richieste in attesa
const processQueue = (error: Error | null, token: string | null = null) => {
    failedQueue.forEach(prom => {
        error ? prom.reject(error) : prom.resolve(token);
    });
    failedQueue = [];
};

// Response Interceptor
axiosInstance.interceptors.response.use(
    (response) => response,
    async (error: AxiosError) => {
        const originalRequest = error.config as InternalAxiosRequestConfig & { 
            _retry?: boolean 
        };
        
        if (error.response?.status === 401 && !originalRequest._retry) {
            // Se sessione morta, rifiuta immediatamente
            if (isSessionDead) return Promise.reject(error);
            
            // CASO 1: Refresh gia' in corso -> metti in coda
            if (isRefreshing) {
                return new Promise((resolve, reject) => {
                    failedQueue.push({ resolve, reject });
                }).then(() => axiosInstance(originalRequest));
            }
            
            // CASO 2: Primo refresh -> acquisisci lock
            originalRequest._retry = true;
            isRefreshing = true;
            
            try {
                await axiosInstance.post('/auth/refresh');
                processQueue(null, null);
                return axiosInstance(originalRequest);
            } catch (refreshError) {
                processQueue(refreshError as Error, null);
                isSessionDead = true;
                emitToast.error('Sessione scaduta. Effettua login.');
                window.dispatchEvent(new CustomEvent('auth:sessionExpired'));
                return Promise.reject(refreshError);
            } finally {
                isRefreshing = false;
            }
        }
        
        return Promise.reject(error);
    }
);
\end{lstlisting}

\subsection{Componente DeckCard}\label{subsec:deckcard-component}

Il componente DeckCard e' ottimizzato con React.memo.

\begin{lstlisting}[language=TypeScript, caption={DeckCard Component}]
export interface DeckCardProps {
    deck: Deck;
    onStudy: (deckId: string) => void;
    onMagicGenerate: (deck: Deck) => void;
    onAddCard: (deckId: string) => void;
    onViewDetail: (deckId: string) => void;
    onDelete: (deck: Deck) => void;
    onDragStart?: () => void;
    onDragEnd?: () => void;
    onTogglePin?: (deck: Deck) => void;
}

const DeckCardComponent: React.FC<DeckCardProps> = ({
    deck, onStudy, onMagicGenerate, onAddCard,
    onViewDetail, onDelete, onTogglePin,
    onDragStart: onDragStartProp, onDragEnd: onDragEndProp
}) => {
    const [showMenu, setShowMenu] = useState(false);
    const [isDragging, setIsDragging] = useState(false);
    
    // Computed values
    const hasDueCards = (deck.dueCount ?? 0) > 0;
    const totalCards = deck.totalCards ?? deck.cards?.length ?? 0;
    const masteredCards = deck.cards?.filter(
        c => c.status === 'mastered'
    ).length ?? 0;
    const masteryPercent = totalCards > 0 
        ? Math.round((masteredCards / totalCards) * 100) 
        : 0;
    
    // Theme calcolato dal titolo (consistente)
    const theme = useMemo(() => getDeckTheme(deck), [deck.title]);
    
    // Drag & Drop
    const handleDragStart = useCallback((e: React.DragEvent) => {
        e.dataTransfer.effectAllowed = 'move';
        e.dataTransfer.setData('deckId', deck.id);
        
        // Anteprima personalizzata
        const dragPreview = document.createElement('div');
        dragPreview.innerHTML = `
            <div style="padding: 12px; background: rgba(139, 92, 246, 0.95);">
                <span>${deck.title}</span>
            </div>
        `;
        document.body.appendChild(dragPreview);
        e.dataTransfer.setDragImage(dragPreview, 100, 40);
        
        setTimeout(() => document.body.removeChild(dragPreview), 0);
        setIsDragging(true);
        onDragStartProp?.();
    }, [deck.id, deck.title, onDragStartProp]);
    
    return (
        <motion.div
            animate={{
                scale: isDragging ? 0.95 : 1,
                opacity: isDragging ? 0.6 : 1,
                rotateZ: isDragging ? 2 : 0
            }}
            draggable
            onDragStart={handleDragStart}
            onDragEnd={handleDragEnd}
            className={...}
        >
            {/* Badges */}
            {deck.pinned && (
                <motion.div className="badge-pinned">⭐</motion.div>
            )}
            {hasDueCards && (
                <motion.div className="badge-due">
                    <Clock className="w-3 h-3" />
                    {deck.dueCount} da ripassare
                </motion.div>
            )}
            
            {/* Content */}
            <DeckCardHeader deck={deck} theme={theme} />
            <DeckCardProgress masteryPercent={masteryPercent} />
            <DeckCardActions 
                deck={deck} 
                onStudy={onStudy}
                onMagicGenerate={onMagicGenerate}
            />
        </motion.div>
    );
};

// Memoizzazione per prevenire re-render
export const DeckCard = React.memo(DeckCardComponent);
\end{lstlisting}

\subsection{Flusso Sessione di Studio}\label{subsec:study-session}

\begin{lstlisting}[language=TypeScript, caption={Study Session Flow}]
// 1. Caricamento sessione
const loadSession = async () => {
    const session = await studyService.getSession(deckId, {
        mode: selectedMode,      // 'flashcard' | 'quiz' | 'typing'
        focus: selectedFocus,    // 'smart' | 'due' | 'weak' | 'all'
        limit: cardLimit,        // numero max cards
        direction: cardDirection // 'front' | 'back' | 'mixed'
    });
    setCards(session.cards);
    setCurrentIndex(0);
};

// 2. Review di una card
const handleReview = async (rating: ReviewRating) => {
    // rating: 1 (Again), 3 (Good), 5 (Easy)
    const result = await studyService.submitReview(deckId, {
        cardId: currentCard.id,
        rating
    });
    
    // Aggiorna stato locale
    updateCardInState(result.card);
    
    // Prossima card o completamento
    if (currentIndex < cards.length - 1) {
        setCurrentIndex(i => i + 1);
    } else {
        completeSession();
    }
};

// 3. Completamento sessione
const completeSession = async () => {
    await studyService.completeSession(deckId, {
        mode: selectedMode,
        stats: {
            correct: correctCount,
            wrong: wrongCount,
            timeSeconds: elapsedTime
        }
    });
    navigate('/study');
};
\end{lstlisting}

% ============================================
% SEZIONE 4: COMUNICAZIONE
% ============================================
\section{Flussi di Comunicazione}\label{sec:communication}

\subsection{Sequence: Review Card}

Il flusso completo di una review:

\begin{enumerate}
    \item \textbf{User} clicca rating (1-5) sulla Flashcard UI
    \item \textbf{StudySessionPage} chiama studyService.submitReview(deckId, payload)
    \item \textbf{studyService} invia POST /api/study/:id/review con CSRF token
    \item \textbf{studyController} valida input e chiama studyService.processCardReview()
    \item \textbf{StudyService} recupera deck da MongoDB con filter \{user: userId\}
    \item \textbf{StudyService} chiama AlgorithmFactory.processReview(card, rating, 'sm2')
    \item \textbf{SM2Algorithm} calcola nuovi parametri (EF', interval, nextReviewDate)
    \item \textbf{StudyService} salva deck aggiornato con await deck.save()
    \item \textbf{StudyService} ritorna \{card, stats\} al controller
    \item \textbf{Controller} ritorna JSON \{success: true, data: result\}
    \item \textbf{studyService} normalizza la risposta
    \item \textbf{StudySessionPage} aggiorna stato e mostra prossima card
\end{enumerate}

\subsection{Sequence: Magic Generate}

Il flusso di generazione AI da PDF:

\begin{enumerate}
    \item \textbf{User} seleziona file PDF e clicca "Generate"
    \item \textbf{UI} chiama studyService.generateFromPDF(deckId, file)
    \item \textbf{studyService} crea FormData e invia POST /:id/generate-pdf
    \item \textbf{multer middleware} salva file su disco (uploads/pdfs/)
    \item \textbf{aiLimiter} verifica rate limit (10 chiamate/ora)
    \item \textbf{studyController.uploadAndGenerate} chiama studyService.generateCardsFromPDF()
    \item \textbf{StudyService}:
    \begin{itemize}
        \item Parsa PDF con pdf-parse
        \item Analizza struttura documento con OpenAI
        \item Divide in semantic chunks (12k char)
        \item Genera flashcard in batch (2 chunk/chiamata)
        \item Deduplica con Jaccard similarity
    \end{itemize}
    \item \textbf{StudyService} salva cards nel deck MongoDB
    \item \textbf{Controller} ritorna \{generatedCount, deck\}
    \item \textbf{UI} aggiorna lista cards con nuove flashcard
\end{enumerate}

% ============================================
% SEZIONE 5: SICUREZZA
% ============================================
\section{Sicurezza}\label{sec:security}

\subsection{Authentication Stack}

\begin{lstlisting}[language=JavaScript]
// 1. requireAuth - Verifica JWT
router.use(requireAuth);

// 2. tenantContext - Estrae userId dal JWT
router.use(tenantContext({ required: true }));

// 3. Rate Limiting per operazioni costose
router.post('/:id/generate-pdf', aiLimiter, handler);
router.post('/exam-solver', examSolverLimiter, handler);

// 4. CSRF Protection
app.use('/api', ensureCsrfCookie);
app.use('/api', requireCsrf);
\end{lstlisting}

\subsection{Input Validation}

\begin{lstlisting}[language=JavaScript]
// Card content validation
if (!front?.trim() || front.length > 2000) {
    throw AppError.validation('Front: max 2000 caratteri');
}
if (!back?.trim() || back.length > 4000) {
    throw AppError.validation('Back: max 4000 caratteri');
}

// File upload validation
const maxSize = 10 * 1024 * 1024; // 10MB
if (req.file.size > maxSize) {
    throw AppError.validation('File troppo grande');
}

// PDF magic bytes check
const pdfValidation = await validatePdfFile(filePath);
if (!pdfValidation.isValid) {
    throw AppError.validation(`PDF corrotto: ${pdfValidation.error}`);
}
\end{lstlisting}

\subsection{Data Isolation}

\begin{lstlisting}[language=JavaScript]
// Esempio: Trova deck per ID (con tenant isolation)
const deck = await Deck.findOne({ 
    _id: deckId, 
    user: userId  // CRITICO: Isolamento dati
});

// Se il deck non esiste O non appartiene all'utente -> 404
if (!deck) {
    throw AppError.notFound('Mazzo');
}
\end{lstlisting}

% ============================================
% SEZIONE 6: CONFIGURAZIONE
% ============================================
\section{Configurazione}\label{sec:config}

\begin{longtable}{|p{4.5cm}|p{2.5cm}|p{7cm}|}
\hline
\textbf{Variabile} & \textbf{Default} & \textbf{Descrizione} \\
\hline
\endfirsthead
OPENAI\_API\_KEY & required & API Key per OpenAI \\
\hline
OPENAI\_MODEL & gpt-4o-mini & Modello LLM \\
\hline
MONGO\_URI & required & Connection string MongoDB \\
\hline
JWT\_SECRET & required & Secret per firma JWT \\
\hline
JWT\_EXPIRES\_IN & 15m & Durata access token \\
\hline
JWT\_REFRESH\_EXPIRES\_IN & 7d & Durata refresh token \\
\hline
CSRF\_SECRET & required & Secret per CSRF tokens \\
\hline
RATE\_LIMIT\_AI\_MAX & 10 & Max chiamate AI/ora per utente \\
\hline
\end{longtable}

% ============================================
% APPENDICE
% ============================================
\appendix
\section{Glossario}\label{app:glossary}

\begin{description}[leftmargin=*]
    \item[SM-2] Algorithm di Spaced Repetition sviluppato da Piotr Wozniak.
    \item[FSRS] Free Spaced Repetition Scheduler - Algoritmo ML-based.
    \item[Source Grounding] Tracciamento del testo originale nel PDF.
    \item[Semantic Chunking] Divisione documento in segmenti significativi.
    \item[Multi-Tenancy] Isolamento dati per utente.
    \item[Easiness Factor] Parametro SM-2 (default 2.5).
\end{description}

\end{document}
